\documentclass{amsart}
% For dealing with references we use the comment environment
\usepackage{verbatim}
\newenvironment{reference}{\comment}{\endcomment}
%\newenvironment{reference}{}{}
\newenvironment{slogan}{\comment}{\endcomment}
\newenvironment{history}{\comment}{\endcomment}

% For commutative diagrams we use Xy-pic
\usepackage[all]{xy}

% We use 2cell for 2-commutative diagrams.
\xyoption{2cell}
\UseAllTwocells

% We use multicol for the list of chapters between chapters
\usepackage{multicol}

% This is generally recommended for better output
\usepackage{lmodern}
\usepackage[T1]{fontenc}

% For cross-file-references

% Package for hypertext links:
\usepackage{hyperref}

% For any local file, say "hello.tex" you want to link to please

% Theorem environments.
%
\theoremstyle{plain}
\newtheorem{theorem}[subsection]{Theorem}
\newtheorem{proposition}[subsection]{Proposition}
\newtheorem{lemma}[subsection]{Lemma}

\theoremstyle{definition}
\newtheorem{definition}[subsection]{Definition}
\newtheorem{example}[subsection]{Example}
\newtheorem{exercise}[subsection]{Exercise}
\newtheorem{situation}[subsection]{Situation}

\theoremstyle{remark}
\newtheorem{remark}[subsection]{Remark}
\newtheorem{remarks}[subsection]{Remarks}

\numberwithin{equation}{subsection}

% Macros
%
\def\lim{\mathop{\mathrm{lim}}\nolimits}
\def\colim{\mathop{\mathrm{colim}}\nolimits}
\def\Spec{\mathop{\mathrm{Spec}}}
\def\Hom{\mathop{\mathrm{Hom}}\nolimits}
\def\Ext{\mathop{\mathrm{Ext}}\nolimits}
\def\SheafHom{\mathop{\mathcal{H}\!\mathit{om}}\nolimits}
\def\SheafExt{\mathop{\mathcal{E}\!\mathit{xt}}\nolimits}
\def\Sch{\mathit{Sch}}
\def\Mor{\mathop{\mathrm{Mor}}\nolimits}
\def\Ob{\mathop{\mathrm{Ob}}\nolimits}
\def\Sh{\mathop{\mathit{Sh}}\nolimits}
\def\NL{\mathop{N\!L}\nolimits}
\def\CH{\mathop{\mathrm{CH}}\nolimits}
\def\proetale{{pro\text{-}\acute{e}tale}}
\def\etale{{\acute{e}tale}}
\def\QCoh{\mathit{QCoh}}
\def\Ker{\mathop{\mathrm{Ker}}}
\def\Im{\mathop{\mathrm{Im}}}
\def\Coker{\mathop{\mathrm{Coker}}}
\def\Coim{\mathop{\mathrm{Coim}}}

% Boxtimes
%
\DeclareMathSymbol{\boxtimes}{\mathbin}{AMSa}{"02}

%
% Macros for moduli stacks/spaces
%
\def\QCohstack{\mathcal{QC}\!\mathit{oh}}
\def\Cohstack{\mathcal{C}\!\mathit{oh}}
\def\Spacesstack{\mathcal{S}\!\mathit{paces}}
\def\Quotfunctor{\mathrm{Quot}}
\def\Hilbfunctor{\mathrm{Hilb}}
\def\Curvesstack{\mathcal{C}\!\mathit{urves}}
\def\Polarizedstack{\mathcal{P}\!\mathit{olarized}}
\def\Complexesstack{\mathcal{C}\!\mathit{omplexes}}
% \Pic is the operator that assigns to X its picard group, usage \Pic(X)
% \Picardstack_{X/B} denotes the Picard stack of X over B
% \Picardfunctor_{X/B} denotes the Picard functor of X over B
\def\Pic{\mathop{\mathrm{Pic}}\nolimits}
\def\Picardstack{\mathcal{P}\!\mathit{ic}}
\def\Picardfunctor{\mathrm{Pic}}
\def\Deformationcategory{\mathcal{D}\!\mathit{ef}}

\setlength{\emergencystretch}{3em}
\begin{document}
\title[Stacks Project, Tag 0EYU]{432 --- Smooth base change for prime-to-characteristic torsion etale cohomology}
\maketitle
\noindent Copyright (C) 2005--2025 Johan de Jong. Stacks Project authors. Source: \href{https://stacks.math.columbia.edu/tag/0EYU}{Tag 0EYU}.

\noindent Editorial difficulty note (not author text). Difficulty H3: The proof must efface pulled-back injective sheaf cohomology, despite pullback not preserving injectivity. It reduces to relative curves, approximates by finite-type bases, passes through finite normalizations and controlled Galois covers of the generic fibre, then combines degree induction and dimension induction to kill the remaining support..

\noindent Editorial provenance note (not author text). History: excerpt from revision \texttt{a0444\allowbreak 6e57e\allowbreak c1fbc\allowbreak 252a8\allowbreak 71afc\allowbreak ec775\allowbreak 2fb28\allowbreak 07b14}, etale-cohomology.tex, lines 17445--17968. Reference presentation was adapted; original-author context and core support blocks were retained or added. The mathematical wording is unchanged. Licensed under GNU FDL 1.2 or later; no invariant sections or cover texts. See ../licenses/COPYING. Context blocks reproduce author definitions and supporting results; they are not additional counted problems.

\begin{theorem}[Smooth base change]
\label{theorem-smooth-base-change}
Consider a cartesian diagram of schemes
$$
\xymatrix{
X \ar[d]_f & Y \ar[l]^h \ar[d]^e \\
S & T \ar[l]_g
}
$$
where $f$ is smooth and $g$ quasi-compact and quasi-separated. Then
$$
f^{-1}R^qg_*\mathcal{F} = R^qh_*e^{-1}\mathcal{F}
$$
for any $q$ and any abelian sheaf $\mathcal{F}$
on $T_\etale$ all of whose stalks at geometric points are torsion of
orders invertible on $S$.
\end{theorem}

\begin{proof}[First proof of smooth base change]
This proof is very long but more direct (using less general theory)
than the second proof given below.

\medskip\noindent
The theorem is local on $X_\etale$. More precisely, suppose we have
$U \to X$ \'etale such that $U \to S$ factors as $U \to V \to S$
with $V \to S$ \'etale. Then we can consider the cartesian square
$$
\xymatrix{
U \ar[d]_{f'} & U \times_X Y \ar[l]^{h'} \ar[d]^{e'} \\
V & V \times_S T \ar[l]_{g'}
}
$$
and setting $\mathcal{F}' = \mathcal{F}|_{V \times_S T}$
we have $f^{-1}R^qg_*\mathcal{F}|_U = (f')^{-1}R^qg'_*\mathcal{F}'$
and $R^qh_*e^{-1}\mathcal{F}|_U = R^qh'_*(e')^{-1}\mathcal{F}'$
(as follows from the compatibility of localization with morphisms of sites, see
Sites, Lemma \href{https://stacks.math.columbia.edu/tag/03EF}{lemma localize morphism strong (Tag 03EF)} and
and
Cohomology on Sites, Lemma
\href{https://stacks.math.columbia.edu/tag/0D6G}{lemma restrict direct image open (Tag 0D6G)}).
Thus it suffices to produce an \'etale covering of $X$ by
$U \to X$ and factorizations $U \to V \to S$
as above such that the theorem holds for the diagram with
$f'$, $h'$, $g'$, $e'$.

\medskip\noindent
By the local structure of smooth morphisms, see
Morphisms, Lemma \href{https://stacks.math.columbia.edu/tag/054L}{lemma smooth etale over affine space (Tag 054L)},
we may assume $X$ and $S$ are affine and $X \to S$
factors through an \'etale morphism $X \to \mathbf{A}^d_S$.
If we have a tower of cartesian diagrams
$$
\xymatrix{
W \ar[d]_i & Z \ar[l]^j \ar[d]^k \\
X \ar[d]_f & Y \ar[l]^h \ar[d]^e \\
S & T \ar[l]_g
}
$$
and the theorem holds for the bottom and top squares, then
the theorem holds for the outer rectangle; this is formal.
Writing $X \to S$ as the composition
$$
X \to \mathbf{A}^{d - 1}_S \to \mathbf{A}^{d - 2}_S \to \ldots \to
\mathbf{A}^1_S \to S
$$
we conclude that it suffices to prove the theorem when $X$ and $S$
are affine and $X \to S$ has relative dimension $1$.

\medskip\noindent
For every $n \geq 1$ invertible on $S$, let $\mathcal{F}[n]$
be the subsheaf of sections of $\mathcal{F}$ annihilated by $n$. Then
$\mathcal{F} = \colim \mathcal{F}[n]$ by our assumption on
the stalks of $\mathcal{F}$. The functors $e^{-1}$ and $f^{-1}$
commute with colimits as they are left adjoints. The functors
$R^qh_*$ and $R^qg_*$ commute with filtered colimits by
Lemma \href{https://stacks.math.columbia.edu/tag/09Z1}{lemma relative colimit (Tag 09Z1)}.
Thus it suffices to prove the theorem for $\mathcal{F}[n]$.
From now on we fix an integer $n$, we work with
sheaves of $\mathbf{Z}/n\mathbf{Z}$-modules and
we assume $S$ is a scheme over $\Spec(\mathbf{Z}[1/n])$.

\medskip\noindent
Next, we reduce to the case where $T$ is affine. Since $g$ is quasi-compact
and quasi-separate and $S$ is affine, the scheme $T$ is quasi-compact and
quasi-separated. Thus we can use the induction principle of
Cohomology of Schemes, Lemma \href{https://stacks.math.columbia.edu/tag/08DR}{lemma induction principle (Tag 08DR)}.
Hence it suffices to show that if $T = W \cup W'$
is an open covering and the theorem holds for the squares
$$
\xymatrix{
X \ar[d] & e^{-1}(W) \ar[l]^i \ar[d] \\
S & W \ar[l]_a
}
\quad
\xymatrix{
X \ar[d] & e^{-1}(W') \ar[l]^j \ar[d] \\
S & W' \ar[l]_b
}
\quad
\xymatrix{
X \ar[d] & e^{-1}(W \cap W') \ar[l]^-k \ar[d] \\
S & W \cap W' \ar[l]_c
}
$$
then the theorem holds for the original diagram. To see this we consider the
diagram
$$
\xymatrix{
f^{-1}R^{q - 1}c_*\mathcal{F}|_{W \cap W'} \ar[d]^{\cong} \ar[r] &
f^{-1}R^qg_*\mathcal{F} \ar[d] \ar[r] &
f^{-1}R^qa_*\mathcal{F}|_W \oplus
f^{-1}R^qb_*\mathcal{F}|_{W'} \ar[d]_{\cong} \\
R^qk_*e^{-1}\mathcal{F}|_{e^{-1}(W \cap W')} \ar[r] &
R^qh_*e^{-1}\mathcal{F} \ar[r] &
R^qi_*e^{-1}\mathcal{F}|_{e^{-1}(W)} \oplus
R^qj_*e^{-1}\mathcal{F}|_{e^{-1}(W')}
}
$$
whose rows are the long exact sequences of
Lemma \href{https://stacks.math.columbia.edu/tag/0EYK}{lemma relative mayer vietoris (Tag 0EYK)}.
Thus the $5$-lemma gives the desired conclusion.

\medskip\noindent
Summarizing, we may assume $S$, $X$, $T$, and $Y$ affine,
$\mathcal{F}$ is $n$ torsion, $X \to S$ is smooth of relative dimension $1$,
and $S$ is a scheme over $\mathbf{Z}[1/n]$.
We will prove the theorem by induction on $q$. The base case $q = 0$
is handled by Lemma \href{https://stacks.math.columbia.edu/tag/0EYS}{lemma fppf reduced fibres base change f star (Tag 0EYS)}.
Assume $q > 0$ and the theorem holds for all smaller degrees.
Choose a short exact sequence
$0 \to \mathcal{F} \to \mathcal{I} \to \mathcal{Q} \to 0$
where $\mathcal{I}$ is an injective sheaf of $\mathbf{Z}/n\mathbf{Z}$-modules.
Consider the induced diagram
$$
\xymatrix{
f^{-1}R^{q - 1}g_*\mathcal{I} \ar[d]_{\cong} \ar[r] &
f^{-1}R^{q - 1}g_*\mathcal{Q} \ar[d]_{\cong} \ar[r] &
f^{-1}R^qg_*\mathcal{F} \ar[d] \ar[r] &
0 \ar[d] \\
R^{q - 1}h_*e^{-1}\mathcal{I} \ar[r] &
R^{q - 1}h_*e^{-1}\mathcal{Q} \ar[r] &
R^qh_*e^{-1}\mathcal{F} \ar[r] &
R^qh_*e^{-1}\mathcal{I}
}
$$
with exact rows. We have the zero in the right upper corner
as $\mathcal{I}$ is injective. The left two vertical arrows are
isomorphisms by induction hypothesis. Thus it suffices to prove
that $R^qh_*e^{-1}\mathcal{I} = 0$.

\medskip\noindent
Write $S = \Spec(A)$ and $T = \Spec(B)$ and say the morphism
$T \to S$ is given by the ring map $A \to B$. We can write
$A \to B = \colim_{i \in I} (A_i \to B_i)$ as a filtered
colimit of maps of rings of finite type over $\mathbf{Z}[1/n]$
(see Algebra, Lemma \href{https://stacks.math.columbia.edu/tag/00QY}{lemma limit no condition (Tag 00QY)}).
For $i \in I$ we set $S_i = \Spec(A_i)$ and $T_i = \Spec(B_i)$.
For $i$ large enough we can find a smooth morphism $X_i \to S_i$
of relative dimension $1$ such that $X = X_i \times_{S_i} S$, see
Limits, Lemmas \href{https://stacks.math.columbia.edu/tag/01ZM}{lemma descend finite presentation (Tag 01ZM)},
\href{https://stacks.math.columbia.edu/tag/0C0C}{lemma descend smooth (Tag 0C0C)}, and
\href{https://stacks.math.columbia.edu/tag/0EY2}{lemma descend dimension d (Tag 0EY2)}. Set $Y_i = X_i \times_{S_i} T_i$
to get squares
$$
\xymatrix{
X_i \ar[d]_{f_i} & Y_i \ar[l]^{h_i} \ar[d]^{e_i} \\
S_i & T_i \ar[l]_{g_i}
}
$$
Observe that $\mathcal{I}_i = (T \to T_i)_*\mathcal{I}$
is an injective sheaf of $\mathbf{Z}/n\mathbf{Z}$-modules on
$T_i$, see Cohomology on Sites, Lemma
\href{https://stacks.math.columbia.edu/tag/0730}{lemma pushforward injective flat (Tag 0730)}.
We have $\mathcal{I} = \colim (T \to T_i)^{-1}\mathcal{I}_i$
by Lemma \href{https://stacks.math.columbia.edu/tag/0EYN}{lemma linus hamann (Tag 0EYN)}. Pulling back by $e$ we get
$e^{-1}\mathcal{I} = \colim (Y \to Y_i)^{-1}e_i^{-1}\mathcal{I}_i$.
By Lemma \href{https://stacks.math.columbia.edu/tag/0EYM}{lemma relative colimit general (Tag 0EYM)} applied to the system
of morphisms $Y_i \to X_i$ with limit $Y \to X$ we have
$$
R^qh_*e^{-1}\mathcal{I} =
\colim (X \to X_i)^{-1} R^qh_{i, *} e_i^{-1}\mathcal{I}_i
$$
This reduces us to the case where $T$ and $S$ are affine of finite
type over $\mathbf{Z}[1/n]$.

\medskip\noindent
Summarizing, we have an integer $q \geq 1$ such that the theorem holds
in degrees $< q$, the schemes $S$ and $T$ affine of finite type
type over $\mathbf{Z}[1/n]$, we have
$X \to S$ smooth of relative dimension $1$ with $X$ affine, and
$\mathcal{I}$ is an injective sheaf of $\mathbf{Z}/n\mathbf{Z}$-modules
and we have to show that $R^qh_*e^{-1}\mathcal{I} = 0$.
We will do this by induction on $\dim(T)$.

\medskip\noindent
The base case is $T = \emptyset$, i.e., $\dim(T) < 0$.
If you don't like this, you can take as your base case
the case $\dim(T) = 0$. In this case $T \to S$ is finite
(in fact even $T \to \Spec(\mathbf{Z}[1/n])$ is finite as the
target is Jacobson; details omitted), so
$h$ is finite too and hence has
vanishing higher direct images (see references below).

\medskip\noindent
Assume $\dim(T) = d \geq 0$ and we know the result for all situations where $T$
has lower dimension. Pick $U$ affine and \'etale over $X$ and a section $\xi$
of $R^qh_*q^{-1}\mathcal{I}$ over $U$. We have to show
that $\xi$ is zero. Of course, we may replace $X$ by $U$
(and correspondingly $Y$ by $U \times_X Y$)
and assume $\xi \in H^0(X, R^qh_*e^{-1}\mathcal{I})$.
Moreover, since $R^qh_*e^{-1}\mathcal{I}$ is a sheaf,
it suffices to prove that $\xi$ is zero locally on $X$.
Hence we may replace $X$ by the members of an \'etale covering.
In particular, using Lemma \href{https://stacks.math.columbia.edu/tag/03Q8}{lemma higher direct images (Tag 03Q8)}
we may assume that $\xi$ is the image of an element
$\tilde \xi \in H^q(Y, e^{-1}\mathcal{I})$.
In terms of $\tilde \xi$ our task is to show that
$\tilde \xi$ dies in $H^q(U_i \times_X Y, e^{-1}\mathcal{I})$
for some \'etale covering $\{U_i \to X\}$.

\medskip\noindent
By More on Morphisms, Lemma \href{https://stacks.math.columbia.edu/tag/0EY4}{lemma cover smooth by special (Tag 0EY4)}
we may assume that $X \to S$ factors as $X \to V \to S$
where $V \to S$ is \'etale and $X \to V$ is a smooth morphism
of affine schemes of relative dimension $1$, has a section, and
has geometrically connected fibres. Observe that
$\dim(V \times_S T) \leq \dim(T) = d$
for example by More on Algebra, Lemma
\href{https://stacks.math.columbia.edu/tag/07QP}{lemma dimension etale extension (Tag 07QP)}.
Hence we may then replace $S$ by $V$ and $T$ by $V \times_S T$
(exactly as in the discussion in the first paragraph of the proof).
Thus we may assume $X \to S$ is smooth of relative dimension $1$,
geometrically connected fibres, and has a section $\sigma : S \to X$.

\medskip\noindent
Let $\pi : T' \to T$ be a finite surjective morphism.
We will use below that $\dim(T') \leq \dim(T) = d$, see
Algebra, Lemma \href{https://stacks.math.columbia.edu/tag/00OJ}{lemma integral dim up (Tag 00OJ)}.
Choose an injective map $\pi^{-1}\mathcal{I} \to \mathcal{I}'$
into an injective sheaf of $\mathbf{Z}/n\mathbf{Z}$-modules.
Then $\mathcal{I} \to \pi_*\mathcal{I}'$ is injective
and hence has a splitting (as $\mathcal{I}$ is an injective
sheaf of $\mathbf{Z}/n\mathbf{Z}$-modules).
Denote $\pi' : Y' = Y \times_T T' \to Y$ the base change of $\pi$
and $e' : Y' \to T'$ the base change of $e$. Picture
$$
\xymatrix{
X \ar[d]_f & Y \ar[l]^h \ar[d]^e & Y' \ar[l]^{\pi'} \ar[d]^{e'} \\
S & T \ar[l]_g & T' \ar[l]_\pi
}
$$
By Proposition \href{https://stacks.math.columbia.edu/tag/03QP}{proposition finite higher direct image zero (Tag 03QP)} and
Lemma \href{https://stacks.math.columbia.edu/tag/0959}{lemma finite pushforward commutes with base change (Tag 0959)} we have
$R\pi'_*(e')^{-1}\mathcal{I}' = e^{-1}\pi_*\mathcal{I}'$.
Thus by the Leray spectral sequence
(Cohomology on Sites, Lemma \href{https://stacks.math.columbia.edu/tag/0732}{lemma Leray (Tag 0732)})
we have
$$
H^q(Y', (e')^{-1}\mathcal{I}') =
H^q(Y, e^{-1}\pi_*\mathcal{I}') \supset H^q(Y, e^{-1}\mathcal{I})
$$
and this remains true after base change by any $U \to X$ \'etale.
Thus we may replace $T$ by $T'$, $\mathcal{I}$ by $\mathcal{I}'$
and $\tilde \xi$ by its image in $H^q(Y', (e')^{-1}\mathcal{I}')$.

\medskip\noindent
Suppose we have a factorization $T \to S' \to S$ where
$\pi : S' \to S$ is finite.
Setting $X' = S' \times_S X$ we can consider the induced diagram
$$
\xymatrix{
X \ar[d]_f & X' \ar[l]^{\pi'} \ar[d]^{f'} & Y \ar[l]^{h'} \ar[d]^e \\
S & S' \ar[l]_\pi & T \ar[l]_g
}
$$
Since $\pi'$ has vanishing higher direct images we see that
$R^qh_*e^{-1}\mathcal{I} = \pi'_*R^qh'_*e^{-1}\mathcal{I}$
by the Leray spectral sequence. Hence
$H^0(X, R^qh_*e^{-1}\mathcal{I}) = H^0(X', R^qh'_*e^{-1}\mathcal{I})$.
Thus $\xi$ is zero if and only if the corresponding section of
$R^qh'_*e^{-1}\mathcal{I}$ is zero\footnote{This
step can also be seen another way. Namely, we have to show that
there is an \'etale covering $\{U_i \to X\}$ such that
$\tilde \xi$ dies in $H^q(U_i \times_X Y, e^{-1}\mathcal{I})$.
However, if we prove there is an \'etale covering
$\{U'_j \to X'\}$ such that $\tilde \xi$ dies in
$H^q(U'_i \times_{X'} Y, e^{-1}\mathcal{I})$, then by property (B)
for $X' \to X$ (Lemma \href{https://stacks.math.columbia.edu/tag/04DQ}{lemma finite B (Tag 04DQ)}) there exists an \'etale covering
$\{U_i \to X\}$ such that $U_i \times_X X'$ is a disjoint union
of schemes over $X'$ each of which factors through $U'_j$ for some $j$.
Thus we see that $\tilde \xi$ dies in $H^q(U_i \times_X Y, e^{-1}\mathcal{I})$
as desired.}.
Thus we may replace $S$ by $S'$ and $X$ by $X'$.
Observe that $\sigma : S \to X$ base changes to $\sigma' : S' \to X'$
and hence after this replacement it is still true that $X \to S$
has a section $\sigma$ and geometrically connected fibres.

\medskip\noindent
We will use that $S$ and $T$ are Nagata schemes, see
Algebra, Proposition \href{https://stacks.math.columbia.edu/tag/0335}{proposition ubiquity nagata (Tag 0335)}
which will guarantee
that various normalizations are finite, see
Morphisms, Lemmas \href{https://stacks.math.columbia.edu/tag/03GR}{lemma nagata normalization finite (Tag 03GR)} and
\href{https://stacks.math.columbia.edu/tag/035S}{lemma nagata normalization (Tag 035S)}.
In particular, we may first replace $T$ by its normalization
and then replace $S$ by the normalization of $S$ in $T$.
Then $T \to S$ is a disjoint union of dominant
morphisms of integral normal schemes, see
Morphisms, Lemma \href{https://stacks.math.columbia.edu/tag/035L}{lemma normal normalization (Tag 035L)}.
Clearly we may argue one connected component
at a time, hence we may assume $T \to S$ is a
dominant morphism of integral normal schemes.

\medskip\noindent
Let $s \in S$ and $t \in T$ be the generic points. By
Lemma \href{https://stacks.math.columbia.edu/tag/0EYT}{lemma smooth base change fields (Tag 0EYT)} there exist finite field
extensions $K/\kappa(t)$ and $k/\kappa(s)$ such that $k$ is contained in $K$
and a finite \'etale Galois covering $Z \to X_k$ with Galois group $G$
of order dividing a power of $n$ split over $\sigma(\Spec(k))$
such that $\tilde \xi$ maps to zero in $H^q(Z_K, e^{-1}\mathcal{I}|_{Z_K})$.
Let $T' \to T$ be the normalization of $T$ in $\Spec(K)$
and let $S' \to S$ be the normalization of $S$ in $\Spec(k)$.
Then we obtain a commutative diagram
$$
\xymatrix{
S' \ar[d] & T' \ar[l] \ar[d] \\
S & T \ar[l]
}
$$
whose vertical arrows are finite. By the arguments given above we
may and do replace $S$ and $T$ by $S'$ and $T'$ (and correspondingly
$X$ by $X \times_S S'$ and $Y$ by $Y \times_T T'$). After this replacement
we conclude we have a finite \'etale Galois covering
$Z \to X_s$ of the generic fibre of $X \to S$
with Galois group $G$ of order dividing a power of $n$
split over $\sigma(s)$ such that $\tilde \xi$ maps to zero in
$H^q(Z_t, (Z_t \to Y)^{-1}e^{-1}\mathcal{I})$.
Here $Z_t = Z \times_S t = Z \times_s t = Z \times_{X_s} Y_t$.
Since $n$ is invertible on $S$,
by Fundamental Groups, Lemma \href{https://stacks.math.columbia.edu/tag/0EYJ}{lemma extend covering (Tag 0EYJ)}
we can find a finite \'etale morphism $U \to X$ whose restriction to $X_s$
is $Z$.

\medskip\noindent
At this point we replace $X$ by $U$ and $Y$ by $U \times_X Y$.
After this replacement it may
no longer be the case that the fibres of $X \to S$ are geometrically
connected (there still is a section but we won't use this), but what
we gain is that after this replacement $\tilde \xi$ maps to zero
in $H^q(Y_t, e^{-1}\mathcal{I})$, i.e., $\tilde \xi$ restricts to
zero on the generic fibre of $Y \to T$.

\medskip\noindent
Recall that $t$ is the spectrum of the function field of $T$, i.e.,
as a scheme $t$ is the limit of the nonempty affine open subschemes of $T$.
By Lemma \href{https://stacks.math.columbia.edu/tag/03Q6}{lemma directed colimit cohomology (Tag 03Q6)} we conclude there exists
a nonempty open subscheme $V \subset T$ such that $\tilde \xi$ maps to zero in
$H^q(Y \times_T V, e^{-1}\mathcal{I}|_{Y \times_T V})$.

\medskip\noindent
Denote $Z = T \setminus V$. Consider the diagram
$$
\xymatrix{
Y \times_T Z \ar[d]_{e_Z} \ar[r]_{i'} &
Y \ar[d]_e &
Y \times_T V \ar[l]^{j'} \ar[d]^{e_V} \\
Z \ar[r]^i &
T &
V \ar[l]_j
}
$$
Choose an injection $i^{-1}\mathcal{I} \to \mathcal{I}'$
into an injective sheaf of $\mathbf{Z}/n\mathbf{Z}$-modules on $Z$.
Looking at stalks we see that the map
$$
\mathcal{I} \to j_*\mathcal{I}|_V \oplus i_*\mathcal{I}'
$$
is injective and hence splits as $\mathcal{I}$ is an injective sheaf
of $\mathbf{Z}/n\mathbf{Z}$-modules. Thus it suffices to show that
$\tilde \xi$ maps to zero in
$$
H^q(Y, e^{-1}j_*\mathcal{I}|_V) \oplus
H^q(Y, e^{-1}i_*\mathcal{I}')
$$
at least after replacing $X$ by the members of an \'etale covering.
Observe that
$$
e^{-1}j_*\mathcal{I}|_V = j'_*e_V^{-1}\mathcal{I}|_V,\quad
e^{-1}i_*\mathcal{I}' = i'_*e_Z^{-1}\mathcal{I}'
$$
By induction hypothesis on $q$ we see that
$$
R^aj'_*e_V^{-1}\mathcal{I}|_V = 0, \quad a = 1, \ldots, q - 1
$$
By the Leray spectral sequence for $j'$ and the vanishing
above it follows that
$$
H^q(Y, j'_*(e_V^{-1}\mathcal{I}|_V))
\longrightarrow
H^q(Y \times_T V, e_V^{-1}\mathcal{I}_V) =
H^q(Y \times_T V, e^{-1}\mathcal{I}|_{Y \times_T V})
$$
is injective. Thus the vanishing of the image of $\tilde \xi$
in the first summand above because we know $\tilde \xi$ vanishes
in $H^q(Y \times_T V, e^{-1}\mathcal{I}|_{Y \times_T V})$.
Since $\dim(Z) < \dim(T) = d$ by induction the image of $\tilde \xi$
in the second summand
$$
H^q(Y, e^{-1}i_*\mathcal{I}') =
H^q(Y, i'_*e_Z^{-1}\mathcal{I}') =
H^q(Y \times_T Z, e_Z^{-1}\mathcal{I}')
$$
dies after replacing $X$ by the members of a suitable \'etale covering.
This finishes the proof of the smooth base change theorem.
\end{proof}

\begin{proof}[Second proof of smooth base change]
This proof is the same as the longer first proof;
it is shorter only in that we have split out the
arguments used in a number of lemmas.

\medskip\noindent
The case of $q = 0$ is
Lemma \href{https://stacks.math.columbia.edu/tag/0EYS}{lemma fppf reduced fibres base change f star (Tag 0EYS)}.
Thus we may assume $q > 0$ and the result is true
for all smaller degrees.

\medskip\noindent
For every $n \geq 1$ invertible on $S$, let $\mathcal{F}[n]$
be the subsheaf of sections of $\mathcal{F}$ annihilated by $n$. Then
$\mathcal{F} = \colim \mathcal{F}[n]$ by our assumption on
the stalks of $\mathcal{F}$. The functors $e^{-1}$ and $f^{-1}$
commute with colimits as they are left adjoints. The functors
$R^qh_*$ and $R^qg_*$ commute with filtered colimits by
Lemma \href{https://stacks.math.columbia.edu/tag/09Z1}{lemma relative colimit (Tag 09Z1)}.
Thus it suffices to prove the theorem for $\mathcal{F}[n]$.
From now on we fix an integer $n$ invertible on $S$ and we work with
sheaves of $\mathbf{Z}/n\mathbf{Z}$-modules.

\medskip\noindent
By Lemma \href{https://stacks.math.columbia.edu/tag/0EZR}{lemma base change local (Tag 0EZR)} the question is \'etale local
on $X$ and $S$. By the local structure of smooth morphisms, see
Morphisms, Lemma \href{https://stacks.math.columbia.edu/tag/054L}{lemma smooth etale over affine space (Tag 054L)},
we may assume $X$ and $S$ are affine and $X \to S$
factors through an \'etale morphism $X \to \mathbf{A}^d_S$.
Writing $X \to S$ as the composition
$$
X \to \mathbf{A}^{d - 1}_S \to \mathbf{A}^{d - 2}_S \to \ldots \to
\mathbf{A}^1_S \to S
$$
we conclude from Lemma \href{https://stacks.math.columbia.edu/tag/0EZS}{lemma base change compose (Tag 0EZS)}
that it suffices to prove the theorem when $X$ and $S$
are affine and $X \to S$ has relative dimension $1$.

\medskip\noindent
By Lemma \href{https://stacks.math.columbia.edu/tag/0F08}{lemma base change does not hold (Tag 0F08)} it suffices to show that
$R^qh_*\mathbf{Z}/d\mathbf{Z} = 0$ for $d | n$ whenever we have a
cartesian diagram
$$
\xymatrix{
X \ar[d] & Y \ar[d] \ar[l]^h \\
S & \Spec(K) \ar[l]
}
$$
where $X \to S$ is affine and smooth of relative dimension $1$,
$S$ is the spectrum of a normal
domain $A$ with algebraically closed fraction field $L$,
and $K/L$ is an extension of algebraically closed fields.

\medskip\noindent
Recall that $R^qh_*\mathbf{Z}/d\mathbf{Z}$ is the sheaf associated
to the presheaf
$$
U \longmapsto
H^q(U \times_X Y, \mathbf{Z}/d\mathbf{Z}) =
H^q(U \times_S \Spec(K), \mathbf{Z}/d\mathbf{Z})
$$
on $X_\etale$ (Lemma \href{https://stacks.math.columbia.edu/tag/03Q8}{lemma higher direct images (Tag 03Q8)}).
Thus it suffices to show: given $U$ and
$\xi \in H^q(U \times_S \Spec(K), \mathbf{Z}/d\mathbf{Z})$
there exists an \'etale covering $\{U_i \to U\}$
such that $\xi$ dies in $H^q(U_i \times_S \Spec(K), \mathbf{Z}/d\mathbf{Z})$.

\medskip\noindent
Of course we may take $U$ affine. Then $U \times_S \Spec(K)$
is a (smooth) affine curve over $K$ and hence we have vanishing
for $q > 1$ by Theorem \href{https://stacks.math.columbia.edu/tag/03SC}{theorem vanishing affine curves (Tag 03SC)}.

\medskip\noindent
Final case: $q = 1$. We may replace $U$ by the members of an
\'etale covering as in More on Morphisms, Lemma
\href{https://stacks.math.columbia.edu/tag/0EY4}{lemma cover smooth by special (Tag 0EY4)}.
Then $U \to S$ factors as $U \to V \to S$ where
$U \to V$ has geometrically connected fibres,
$U$, $V$ are affine, $V \to S$ is \'etale, and there is
a section $\sigma : V \to U$. By
Lemma \href{https://stacks.math.columbia.edu/tag/09Z9}{lemma normal scheme with alg closed function field (Tag 09Z9)}
we see that $V$ is isomorphic to a (finite) disjoint union of
(affine) open subschemes of $S$.
Clearly we may replace $S$ by one of these
and $X$ by the corresponding component of $U$.
Thus we may assume $X \to S$ has geometrically connected
fibres, has a section $\sigma$, and
$\xi \in H^1(Y, \mathbf{Z}/d\mathbf{Z})$.
Since $K$ and $L$ are algebraically closed we have
$$
H^1(X_L, \mathbf{Z}/d\mathbf{Z}) =
H^1(Y, \mathbf{Z}/d\mathbf{Z})
$$
See Lemma \href{https://stacks.math.columbia.edu/tag/0A5E}{lemma base change dim 1 separably closed (Tag 0A5E)}.
Thus there is a finite \'etale Galois covering
$Z \to X_L$ with Galois group $G \subset \mathbf{Z}/d\mathbf{Z}$
which annihilates $\xi$.
You can either see this by looking at the statement or
proof of Lemma \href{https://stacks.math.columbia.edu/tag/0EYT}{lemma smooth base change fields (Tag 0EYT)}
or by using directly that $\xi$ corresponds to a
$\mathbf{Z}/d\mathbf{Z}$-torsor over $X_L$.
Finally, by
Fundamental Groups, Lemma \href{https://stacks.math.columbia.edu/tag/0EZJ}{lemma extend covering general (Tag 0EZJ)}
we find a (necessarily surjective)
finite \'etale morphism $X' \to X$
whose restriction to $X_L$ is $Z \to X_L$.
Since $\xi$ dies in $X'_K$ this finishes the proof.
\end{proof}
\end{document}
